\documentclass[11pt]{article}
\usepackage[T1]{fontenc}\usepackage{amsmath,amssymb}\usepackage[margin=1in]{geometry}\usepackage[hidelinks]{hyperref}
\title{Proof Dossier: Ragasa2026A Maximized Normalized Correlation Bound}\author{Project Koios Formal}\date{}
\begin{document}\maketitle
\section{Identity and status}
ID \texttt{BG-RAGASA2026A-THM-013}; status \texttt{proved}. Lean declaration:
\par\noindent {\tiny\path{ProjectKoios.BandGauge.Papers.Ragasa2026A.abs_maximumOperatorCorrelation_div_le_one}}.
\section{Exact statement}
If the centered parent has common power $\kappa>0$, then
\[\left|\frac{b_{\max}}{\kappa}\right|\le1.\]
\section{Source, proof, and boundary}
This \texttt{formalizes} the normalized maximum used in manuscript lines 645--650 at owner
commit
\begin{center}\small\texttt{99a46cb59d0c6f27ada7d709f194996ef9619d9f}.\end{center}
The maximizing
angle turns $b_{\max}$ into an ordinary field correlation, to which the proved weighted
Frobenius bound applies. Parent power and Pauli coordinates are hypotheses, not validated
data. Kernel axioms: \texttt{propext}, \texttt{Classical.choice}, \texttt{Quot.sound};
no admissions. AI-assisted and kernel-checked; no independent expert review recorded.
\end{document}
